\documentclass[10pt,a4paper]{amsart}
\usepackage[margin=19mm]{geometry}
\usepackage{amsmath,amssymb,mathtools}
\usepackage[scr=rsfs,cal=euler]{mathalfa}
\usepackage{xfrac}
\usepackage[unicode,hidelinks,bookmarks=false]{hyperref}
\newcommand{\ie}{i.e.\ }
\newcommand{\cf}{cf.\ }
\newcommand{\resp}{respectively\ }
\newcommand{\Subsection}{\subsection}
\providecommand{\nobreakdash}{\nobreak\hbox{-}}
\setlength{\emergencystretch}{2em}
\multlinegap0pt
\usepackage{amssymb}
\usepackage{tikz}
\usepackage{float}
\newtheorem{theorem}{Theorem}
\title{Interval-Valued Optimization Problems for Strongly LU-E-Invex and Strongly LU-E-Preinvex Functions}
\author{Tauheed, Akhlad Iqbal, Amir Suhail}
\date{Source: arXiv:2602.23704v1 (CC BY 4.0)}
\begin{document}
\maketitle

\noindent\emph{Source and licence.} Interval-Valued Optimization Problems for Strongly LU-E-Invex and Strongly LU-E-Preinvex Functions. Version arXiv:2602.23704v1 (CC BY 4.0), distributed by arXiv under the Creative Commons Attribution 4.0 International licence, http://creativecommons.org/licenses/by/4.0/, as recorded in the arXiv licence metadata for this exact version and shown on the abstract page https://arxiv.org/abs/2602.23704v1. Full licence notice: \texttt{licenses/arxiv-2602.23704-CC-BY-4.0.txt}.\par\smallskip

\noindent\textit{Editorial scope.} Counted unit: the article's theorem theorem (source0.tex lines 1562--1568). The article's complete original proof of it is delivered with it (source0.tex lines 1570--1625, one \texttt{proof} environment of 56 lines). No mathematics is written, repaired or re-derived: the statement and the proof are the article's own lines, in the article's own notation, and no retained line is substituted or re-flowed.  Retained uncounted as the source-local context the proof needs: lines 1518--1520.  Nothing was omitted from any retained span, so the package carries no omission record.  No retained line contains a citation, so the wrapper carries no bibliography and no bibliography entry.  No retained line contains a figure environment, an image-inclusion command or drawing code, so no image is delivered and the package declares no asset dependency.  Editorial formatting only, in addition: an AMS \texttt{amsart} preamble that restates the article's own declarations and macros used by the retained lines, namely the environments \texttt{theorem} on separate counters, together with the standard packages those lines need.  The retained environments print in this document as Theorem 1 in order of appearance, whereas the article prints the same items with the numbering of its own full document.  The complete native source of the article as distributed in the arXiv e-print for this exact version is bundled unchanged as \texttt{sources/195418/source0.tex}.
       \begin{equation}\label{4.2} 
          {X}=\big\{\mathit{\zeta} \in\mathbb{R}^{n}~:~\tilde{h}_{j}(\mathit{\zeta} )\preceq 0,~1\leq j\leq m\big\}.
       \end{equation} 

\begin{theorem}
Let $\tilde{h}_{j}: \mathbb{R}^{n} \rightarrow \mathcal{I}(\mathbb{R})$, for $0 \leq j \leq m$, be SLU$\mathcal{E}$P functions on $\mathbb{R}^{n}$, and let $X$ be the non-empty feasible region defined by \emph{(\ref{4.2})}.Consider the set of optimal solutions 
\[
X_{\text{opt}} = \big\{\, \zeta \in X : \tilde{h}_{0}(\zeta) = \min_{w \in X} \tilde{h}_{0}(w) \big\}
\] and $\mathcal{E}(X_{\text{opt}})\subseteq X_{\text{opt}}$. Then, $X_{\text{opt}}$
is an S$\mathcal{E}$I set with respect to $\Psi$.
\end{theorem}

\begin{proof}
By assumption, every $\tilde{h}_{j}$ for $0 \leq j \leq m$ is SLU$\mathcal{E}$P on $\mathbb{R}^{n}$. Let $\mathit{\zeta}, \mathit{\delta} \in X_{\text{opt}}$ and note that, by definition of optimality, we have
\[
\tilde{h}_{0}(\mathit{\zeta}) = \tilde{h}_{0}(\mathit{\delta}) = \tilde{h}_{0}(\mathcal{E}(\mathit{\zeta})) = \tilde{h}_{0}(\mathcal{E}(\mathit{\delta})) = \tilde{h}^{*},
\quad \text{where} \quad 
\tilde{h}^{*} := \min_{w \in X} \tilde{h}_{0}(w).
\]

Since $X$ is feasible and closed under $\mathcal{E}$, it follows that
$\mathit{\zeta}, \mathit{\delta}, \mathcal{E}(\mathit{\zeta}), \mathcal{E}(\mathit{\delta}) \in X$. By the SLU$\mathcal{E}$P property, for any $\alpha, \lambda \in [0,1]$, we have
\[
\tilde{h}_{0}\big( 
\alpha\, \mathit{\delta} + \mathcal{E}(\mathit{\delta})
+ \lambda\, \Psi(
\alpha\, \mathit{\zeta} + \mathcal{E}(\mathit{\zeta}),
\, \alpha\, \mathit{\delta} + \mathcal{E}(\mathit{\delta})
)
\big)
\preceq
\lambda\, \tilde{h}_{0}(\mathcal{E}(\mathit{\zeta})) + (1-\lambda)\, \tilde{h}_{0}(\mathcal{E}(\mathit{\delta})) 
= \tilde{h}^{*}.
\]
But since $\tilde{h}^{*}$ is the minimum value, the above inequality must hold as equality. Therefore,
\[
\tilde{h}_{0}\big(
\alpha\, \mathit{\delta} + \mathcal{E}(\mathit{\delta}) 
+ \lambda\, \Psi(
\alpha\, \mathit{\zeta} + \mathcal{E}(\mathit{\zeta}),
\, \alpha\, \mathit{\delta} + \mathcal{E}(\mathit{\delta})
)
\big) = \tilde{h}^{*}.
\]

Moreover, for all $1 \leq j \leq m$, the constraint functions satisfy
\[
\tilde{h}_{j}\big(
\alpha\, \mathit{\delta} + \mathcal{E}(\mathit{\delta})
+ \lambda\, \Psi(
\alpha\, \mathit{\zeta} + \mathcal{E}(\mathit{\zeta}),
\, \alpha\, \mathit{\delta} + \mathcal{E}(\mathit{\delta})
)
\big) \preceq 0.
\]

Hence, the constructed point lies in $X$ and achieves the optimal value $\tilde{h}^{*}$ for the objective. Therefore, it is also an element of $X_{\text{opt}}$:
\[
\alpha\, \mathit{\delta} + \mathcal{E}(\mathit{\delta})
+ \lambda\, \Psi(
\alpha\, \mathit{\zeta} + \mathcal{E}(\mathit{\zeta}),
\, \alpha\, \mathit{\delta} + \mathcal{E}(\mathit{\delta})
) \in X_{\text{opt}}.
\]
This proves that $X_{\text{opt}}$ is closed under the same generalized operation that defines an S$\mathcal{E}$I set.

Therefore, $X_{\text{opt}}$ is indeed an S$\mathcal{E}$I set with respect to $\Psi$.
\end{proof} 

\end{document}
